\documentclass[a4paper]{article}
\usepackage[pdftex]{graphicx}
\usepackage[utf8]{inputenc}
\usepackage{enumerate}
\usepackage{icomma}
\usepackage{amssymb}
\usepackage{polynom}
\usepackage{tikz}
\usepackage{href-ul}
\hypersetup{
	colorlinks=true,
	linkcolor=blue,
	urlcolor=blue}
\usepackage{geometry}
\geometry{a4paper, top=15mm, left=15mm, right=15mm, bottom=15mm,
	headsep=10mm, footskip=12mm}

\begin{document}
	\begin{tikzpicture}(10,3)
		\draw[ultra thick](2,0) --(10,0) -- (10,1.3) --(2,1.3) -- (2,0);
		\draw[fill=black](2,0)-- (0,.6) -- (2,.6) -- (2,0);
		\node at (6,.6) {\large solution sheet from \href{http://www.okuyakl.com}{www.okuyakl.com}};
	\end{tikzpicture}
	\vspace{0.5 cm}
	
	\noindent{\bf Task 1. a)}\\
	
	We substitute $x=3$ into the function term and set it equal to zero:
	$$
	\renewcommand{\arraystretch}{1.5}
	\begin{array}{rcl}
		0.2 \cdot 3^3 -1.4 \cdot 3^2 - 1.8 \cdot 3 +12.6&=0 \\
		5.4 - 12.6 -5.4 +12.6 &=0 & (w)
	\end{array}
	$$
	We get a true statement.
	
	\noindent {\bf task 1. b)}\\
	We get further zeros by polynomial division using the already known factor
	$(x-3)$:
	$$
	\renewcommand{\arraystretch}{1.5}
	\begin{array}{rrrrcl}
		0.2 x^3 &-1.4 x^2 &- 1.8 x &+12.6 & :(x-3)=& 0.2x^2-0.8x -4.2 \\
		\underline{-(0.2x^3}& \underline{-0.6 x^2)}\\
		& -0.8 x^2 & -1.8 x \\
		&\underline{-(-0.8x^2}& \underline{+2.4 x)}\\
		& & -4.2 x &+12.6 \\
		& & \underline{-(-4.2x}& \underline{+12.6)}\\
		& & & 0 \\
	\end{array}
	$$
	We examine the quadratic equation obtained using the midnight formula:
	$$
	\renewcommand{\arraystretch}{1.5}
	\begin{array}{rcl}
		x_{2/3}&=&\frac{0.8 \pm \sqrt{0.64 -4\cdot 0.2 \cdot (-4.2)}}{2\cdot 0.2}\\
		\underline{x_2} &=& \underline{7} \\
		\underline{x_3} &=& \underline{-3}
	\end{array}
	$$
	Other roots of the function are $N_2(-3|0)$ and $N_3(7|0)$.
	
	\noindent {\bf task 1. c)}\\
	The intersection with the y-axis is obtained by inserting $x=0$ into the function term:
	$$g(0)=0-0-0+12.6 =12.6 \quad \Rightarrow \quad S_y(0|12.6)$$
	
	\noindent {\bf task 1. d)}\\
	The highest occurring power, here $x^3$, uses the sign of its prefactor to decide where the function goes for large $|x|$. The prefactor is positive and the power is odd. For this reason:
	$$\lim\limits_{x \to \infty} 0.2 x^3 -1.4 x^2 - 1.8 x +12.6 = \infty$$
	$$\lim\limits_{x \to -\infty} 0.2 x^3 -1.4 x^2 - 1.8 x +12.6 = -\infty$$
	
	\noindent{\bf Task 2.}\\
	\begin{itemize}
		\item The y-axis is a pole without a sign change $\Rightarrow \quad x^2 $ in the denominator.
		\item Double zero at $x=1 \quad \Rightarrow \quad (x-1)^2$ in the numerator
		\item numerator degree = denominator degree; horizontal asymptote at $y=0.5 \quad \Rightarrow $ factor $ {1 \over 2}$
	\end{itemize}
	Function term:
	$$f(x)= {(x-1)^2 \over 2 x^2}$$
	
	\newpage
	\noindent{\bf Task 3. a)}\\
	Gaps in definition are zeros of the denominator. Using Vieta's theorem we get:
	$$
	\renewcommand{\arraystretch}{2}
	\begin{array}{rcll}
		x^2+2x-3 & = & 0 & = (x+a)(x+b)\\
		a + b & = & 2 \\
		a \cdot b & = & -3 \\
		\Rightarrow a & = & 3 \\
		b & = & -1 \\
		x^2+2x-3 & = & 0 & = (x+3)(x-1)
	\end{array}
	$$
	So the roots of this term are $x_1=-3$ and $x_2=1$.
	
	\noindent{\bf Task 3. b)}\\
	We guess the first zero of the counter, it is $x=1$. We then carry out the polynomial division:
	$$
	\polyset{style=C,div=:,vars=x}
	\polylongdiv{2x^3-2x^2-8x+8}{x-1}
	$$
	We further factor:
	$$2x^2-8 = 2(x^2-4) = 2(x+2)(x-2)$$
	
	This gives us:
	$$f(x)={2(x+2)(x-2)(x-1) \over (x+3)(x-1)}$$
	
	\noindent{\bf Task 3. c)}\\
	We reduce the term $(x-1)$ from the numerator and denominator and look at the behavior of the function in the vicinity of $x=1$:
	$$\lim\limits_{x\to 1^+} {2x^2-8 \over x+3} = -{6 \over 4} = -1.5 $$
	$$\lim\limits_{x\to 1^-} {2x^2-8 \over x+3} = -{6 \over 4} = -1.5 $$
	The limit exists and the function approaches the same point from both sides $y=-1.5$, so the function $f$ can be continued continuously at its definition gap $x=1$.
	
	\noindent{\bf Task 3. d)}\\
	The position $x=-3$ comes into question as a pole position. We also use the limit here:
	$$\lim\limits_{x\to -3^+} {2x^2-8 \over x+3} = "{10 \over +0}" = + \infty $$
	$$\lim\limits_{x\to -3^-} {2x^2-8 \over x+3} = "{10 \over -0}" = - \infty $$
	\noindent{\bf Task 3. e)}\\
	The pole point found at $x=-3$ is also a vertical asymptote. The numerator degree is one greater than the denominator degree, so we still find a slanted asymptote:
	$$
	\polyset{style=C,div=:,vars=x}
	\polylongdiv{2x^3-2x^2-8x+8}{x^2+2x-3}
	$$
	The slanted asymptote is the term before the remainder:
	$$y=2x-6$$.
	
	\noindent{\bf Task 4. a)}\\
	The definition set is:
	$$D_f =\mathbb{R} \setminus \{-2\}$$
	Horizontal asymptote: $y=5$ \\
	Vertical asymptote: $x=2$
	
\noindent{\bf Task 4. b)}\\
$$
\renewcommand{\arraystretch}{2}
\begin{array}{rcll}
	{1 \over x+2} +5 &=& 7 &| -5\\
	{1 \over x+2} &=& 2 &| \cdot (x+2)\\
	1 &=& 2x +4 & \Rightarrow \quad x= -1.5
\end{array}
$$

\noindent{\bf Task 4. c)}\\
$$g(x)= { 1 \over {0.5\cdot x +2}}+6 $$
Horizontal asymptote: $y=6$ \\
Vertical asymptote: $x=4$

\begin{center}
	\includegraphics[width=7 cm]{../../viecher/eendcomic.pdf}
	
	Here you can go back to the \href{https://www.okuyakl.de/math/m10fubL033/ae033.pdf}{task sheet}
\end{center}

\end{document}